\section{}

\paragraph{1232年}\eventanchor{YUAN-1232-REGNAL-YEAR}{YUAN-1232-REGNAL-YEAR}　蒙古以窝阔台在位第5年作为诏令、赋役与外交文书的纪年框架。账本列出《元史》〈本纪〉及相关志；《元朝秘史》；《元典章》或《通制条格》相关条作为来源线索；该材料可核对时间与制度称谓，但有关执行范围、群体经验、军数或后果，仍须与地方文书、物质遗存及相关政权记录互证。

\paragraph{1233年}\eventanchor{YUAN-1233-REGNAL-YEAR}{YUAN-1233-REGNAL-YEAR}　蒙古以窝阔台在位第6年作为诏令、赋役与外交文书的纪年框架。账本列出《元史》〈本纪〉及相关志；《元朝秘史》；《元典章》或《通制条格》相关条作为来源线索；该材料可核对时间与制度称谓，但有关执行范围、群体经验、军数或后果，仍须与地方文书、物质遗存及相关政权记录互证。

\paragraph{1234年}\eventanchor{YUAN-1234-EVENT-160}{YUAN-1234-EVENT-160}　蒙古与南宋攻灭金朝。账本列出《元史》〈本纪〉及相关志；《元朝秘史》；《元典章》或《通制条格》相关条作为来源线索；该材料可核对时间与制度称谓，但有关执行范围、群体经验、军数或后果，仍须与地方文书、物质遗存及相关政权记录互证。

\paragraph{1234年}\eventanchor{YUAN-1234-REGNAL-YEAR}{YUAN-1234-REGNAL-YEAR}　蒙古以窝阔台在位第7年作为诏令、赋役与外交文书的纪年框架。账本列出《元史》〈本纪〉及相关志；《元朝秘史》；《元典章》或《通制条格》相关条作为来源线索；该材料可核对时间与制度称谓，但有关执行范围、群体经验、军数或后果，仍须与地方文书、物质遗存及相关政权记录互证。

\paragraph{1235年}\eventanchor{YUAN-1235-EVENT-161}{YUAN-1235-EVENT-161}　蒙古贵族会议决定多方向扩张。账本列出《元史》〈本纪〉及相关志；《元朝秘史》；《元典章》或《通制条格》相关条作为来源线索；该材料可核对时间与制度称谓，但有关执行范围、群体经验、军数或后果，仍须与地方文书、物质遗存及相关政权记录互证。

\paragraph{1235年}\eventanchor{YUAN-1235-REGNAL-YEAR}{YUAN-1235-REGNAL-YEAR}　蒙古以窝阔台在位第8年作为诏令、赋役与外交文书的纪年框架。账本列出《元史》〈本纪〉及相关志；《元朝秘史》；《元典章》或《通制条格》相关条作为来源线索；该材料可核对时间与制度称谓，但有关执行范围、群体经验、军数或后果，仍须与地方文书、物质遗存及相关政权记录互证。

\paragraph{1236年}\eventanchor{YUAN-1236-REGNAL-YEAR}{YUAN-1236-REGNAL-YEAR}　蒙古以窝阔台在位第9年作为诏令、赋役与外交文书的纪年框架。账本列出《元史》〈本纪〉及相关志；《元朝秘史》；《元典章》或《通制条格》相关条作为来源线索；该材料可核对时间与制度称谓，但有关执行范围、群体经验、军数或后果，仍须与地方文书、物质遗存及相关政权记录互证。

\paragraph{1237年}\eventanchor{YUAN-1237-REGNAL-YEAR}{YUAN-1237-REGNAL-YEAR}　蒙古以窝阔台在位第10年作为诏令、赋役与外交文书的纪年框架。账本列出《元史》〈本纪〉及相关志；《元朝秘史》；《元典章》或《通制条格》相关条作为来源线索；该材料可核对时间与制度称谓，但有关执行范围、群体经验、军数或后果，仍须与地方文书、物质遗存及相关政权记录互证。

\paragraph{1238年}\eventanchor{YUAN-1238-REGNAL-YEAR}{YUAN-1238-REGNAL-YEAR}　蒙古以窝阔台在位第11年作为诏令、赋役与外交文书的纪年框架。账本列出《元史》〈本纪〉及相关志；《元朝秘史》；《元典章》或《通制条格》相关条作为来源线索；该材料可核对时间与制度称谓，但有关执行范围、群体经验、军数或后果，仍须与地方文书、物质遗存及相关政权记录互证。

\paragraph{1239年}\eventanchor{YUAN-1239-REGNAL-YEAR}{YUAN-1239-REGNAL-YEAR}　蒙古以窝阔台在位第12年作为诏令、赋役与外交文书的纪年框架。账本列出《元史》〈本纪〉及相关志；《元朝秘史》；《元典章》或《通制条格》相关条作为来源线索；该材料可核对时间与制度称谓，但有关执行范围、群体经验、军数或后果，仍须与地方文书、物质遗存及相关政权记录互证。

\paragraph{1240年}\eventanchor{YUAN-1240-REGNAL-YEAR}{YUAN-1240-REGNAL-YEAR}　蒙古以窝阔台在位第13年作为诏令、赋役与外交文书的纪年框架。账本列出《元史》〈本纪〉及相关志；《元朝秘史》；《元典章》或《通制条格》相关条作为来源线索；该材料可核对时间与制度称谓，但有关执行范围、群体经验、军数或后果，仍须与地方文书、物质遗存及相关政权记录互证。

\paragraph{1241年}\eventanchor{YUAN-1241-EVENT-162}{YUAN-1241-EVENT-162}　窝阔台去世，帝国进入摄政与继承协商阶段。账本列出《元史》〈本纪〉及相关志；《元朝秘史》；《元典章》或《通制条格》相关条作为来源线索；该材料可核对时间与制度称谓，但有关执行范围、群体经验、军数或后果，仍须与地方文书、物质遗存及相关政权记录互证。

\paragraph{1241年}\eventanchor{YUAN-1241-REGNAL-YEAR}{YUAN-1241-REGNAL-YEAR}　蒙古以窝阔台在位第14年作为诏令、赋役与外交文书的纪年框架。账本列出《元史》〈本纪〉及相关志；《元朝秘史》；《元典章》或《通制条格》相关条作为来源线索；该材料可核对时间与制度称谓，但有关执行范围、群体经验、军数或后果，仍须与地方文书、物质遗存及相关政权记录互证。

\paragraph{1242年}\eventanchor{YUAN-1242-REGNAL-YEAR}{YUAN-1242-REGNAL-YEAR}　蒙古以乃马真后摄政在位第1年作为诏令、赋役与外交文书的纪年框架。账本列出《元史》〈本纪〉及相关志；《元朝秘史》；《元典章》或《通制条格》相关条作为来源线索；该材料可核对时间与制度称谓，但有关执行范围、群体经验、军数或后果，仍须与地方文书、物质遗存及相关政权记录互证。

\paragraph{1243年}\eventanchor{YUAN-1243-REGNAL-YEAR}{YUAN-1243-REGNAL-YEAR}　蒙古以乃马真后摄政在位第2年作为诏令、赋役与外交文书的纪年框架。账本列出《元史》〈本纪〉及相关志；《元朝秘史》；《元典章》或《通制条格》相关条作为来源线索；该材料可核对时间与制度称谓，但有关执行范围、群体经验、军数或后果，仍须与地方文书、物质遗存及相关政权记录互证。

\paragraph{1244年}\eventanchor{YUAN-1244-REGNAL-YEAR}{YUAN-1244-REGNAL-YEAR}　蒙古以乃马真后摄政在位第3年作为诏令、赋役与外交文书的纪年框架。账本列出《元史》〈本纪〉及相关志；《元朝秘史》；《元典章》或《通制条格》相关条作为来源线索；该材料可核对时间与制度称谓，但有关执行范围、群体经验、军数或后果，仍须与地方文书、物质遗存及相关政权记录互证。

\paragraph{1245年}\eventanchor{YUAN-1245-REGNAL-YEAR}{YUAN-1245-REGNAL-YEAR}　蒙古以乃马真后摄政在位第4年作为诏令、赋役与外交文书的纪年框架。账本列出《元史》〈本纪〉及相关志；《元朝秘史》；《元典章》或《通制条格》相关条作为来源线索；该材料可核对时间与制度称谓，但有关执行范围、群体经验、军数或后果，仍须与地方文书、物质遗存及相关政权记录互证。

\paragraph{1246年}\eventanchor{YUAN-1246-EVENT-163}{YUAN-1246-EVENT-163}　贵由即位为大汗。账本列出《元史》〈本纪〉及相关志；《元朝秘史》；《元典章》或《通制条格》相关条作为来源线索；该材料可核对时间与制度称谓，但有关执行范围、群体经验、军数或后果，仍须与地方文书、物质遗存及相关政权记录互证。

\paragraph{1246年}\eventanchor{YUAN-1246-REGNAL-YEAR}{YUAN-1246-REGNAL-YEAR}　蒙古以贵由在位第1年作为诏令、赋役与外交文书的纪年框架。账本列出《元史》〈本纪〉及相关志；《元朝秘史》；《元典章》或《通制条格》相关条作为来源线索；该材料可核对时间与制度称谓，但有关执行范围、群体经验、军数或后果，仍须与地方文书、物质遗存及相关政权记录互证。

\paragraph{1247年}\eventanchor{YUAN-1247-REGNAL-YEAR}{YUAN-1247-REGNAL-YEAR}　蒙古以贵由在位第2年作为诏令、赋役与外交文书的纪年框架。账本列出《元史》〈本纪〉及相关志；《元朝秘史》；《元典章》或《通制条格》相关条作为来源线索；该材料可核对时间与制度称谓，但有关执行范围、群体经验、军数或后果，仍须与地方文书、物质遗存及相关政权记录互证。

\paragraph{1248年}\eventanchor{YUAN-1248-REGNAL-YEAR}{YUAN-1248-REGNAL-YEAR}　蒙古以贵由在位第3年作为诏令、赋役与外交文书的纪年框架。账本列出《元史》〈本纪〉及相关志；《元朝秘史》；《元典章》或《通制条格》相关条作为来源线索；该材料可核对时间与制度称谓，但有关执行范围、群体经验、军数或后果，仍须与地方文书、物质遗存及相关政权记录互证。

\paragraph{1249年}\eventanchor{YUAN-1249-REGNAL-YEAR}{YUAN-1249-REGNAL-YEAR}　蒙古以海迷失后摄政在位第1年作为诏令、赋役与外交文书的纪年框架。账本列出《元史》〈本纪〉及相关志；《元朝秘史》；《元典章》或《通制条格》相关条作为来源线索；该材料可核对时间与制度称谓，但有关执行范围、群体经验、军数或后果，仍须与地方文书、物质遗存及相关政权记录互证。

\paragraph{1250年}\eventanchor{YUAN-1250-REGNAL-YEAR}{YUAN-1250-REGNAL-YEAR}　蒙古以海迷失后摄政在位第2年作为诏令、赋役与外交文书的纪年框架。账本列出《元史》〈本纪〉及相关志；《元朝秘史》；《元典章》或《通制条格》相关条作为来源线索；该材料可核对时间与制度称谓，但有关执行范围、群体经验、军数或后果，仍须与地方文书、物质遗存及相关政权记录互证。

\paragraph{1251年}\eventanchor{YUAN-1251-EVENT-164}{YUAN-1251-EVENT-164}　蒙哥即大汗位并整肃政治对手。账本列出《元史》〈本纪〉及相关志；《元朝秘史》；《元典章》或《通制条格》相关条作为来源线索；该材料可核对时间与制度称谓，但有关执行范围、群体经验、军数或后果，仍须与地方文书、物质遗存及相关政权记录互证。

\paragraph{1251年}\eventanchor{YUAN-1251-REGNAL-YEAR}{YUAN-1251-REGNAL-YEAR}　蒙古以蒙哥在位第1年作为诏令、赋役与外交文书的纪年框架。账本列出《元史》〈本纪〉及相关志；《元朝秘史》；《元典章》或《通制条格》相关条作为来源线索；该材料可核对时间与制度称谓，但有关执行范围、群体经验、军数或后果，仍须与地方文书、物质遗存及相关政权记录互证。

\paragraph{1252年}\eventanchor{YUAN-1252-REGNAL-YEAR}{YUAN-1252-REGNAL-YEAR}　蒙古以蒙哥在位第2年作为诏令、赋役与外交文书的纪年框架。账本列出《元史》〈本纪〉及相关志；《元朝秘史》；《元典章》或《通制条格》相关条作为来源线索；该材料可核对时间与制度称谓，但有关执行范围、群体经验、军数或后果，仍须与地方文书、物质遗存及相关政权记录互证。

\paragraph{1253年}\eventanchor{YUAN-1253-EVENT-165}{YUAN-1253-EVENT-165}　蒙古军征服大理国。账本列出《元史》〈本纪〉及相关志；《元朝秘史》；《元典章》或《通制条格》相关条作为来源线索；该材料可核对时间与制度称谓，但有关执行范围、群体经验、军数或后果，仍须与地方文书、物质遗存及相关政权记录互证。

\paragraph{1253年}\eventanchor{YUAN-1253-REGNAL-YEAR}{YUAN-1253-REGNAL-YEAR}　蒙古以蒙哥在位第3年作为诏令、赋役与外交文书的纪年框架。账本列出《元史》〈本纪〉及相关志；《元朝秘史》；《元典章》或《通制条格》相关条作为来源线索；该材料可核对时间与制度称谓，但有关执行范围、群体经验、军数或后果，仍须与地方文书、物质遗存及相关政权记录互证。

\paragraph{1254年}\eventanchor{YUAN-1254-REGNAL-YEAR}{YUAN-1254-REGNAL-YEAR}　蒙古以蒙哥在位第4年作为诏令、赋役与外交文书的纪年框架。账本列出《元史》〈本纪〉及相关志；《元朝秘史》；《元典章》或《通制条格》相关条作为来源线索；该材料可核对时间与制度称谓，但有关执行范围、群体经验、军数或后果，仍须与地方文书、物质遗存及相关政权记录互证。

\paragraph{1255年}\eventanchor{YUAN-1255-REGNAL-YEAR}{YUAN-1255-REGNAL-YEAR}　蒙古以蒙哥在位第5年作为诏令、赋役与外交文书的纪年框架。账本列出《元史》〈本纪〉及相关志；《元朝秘史》；《元典章》或《通制条格》相关条作为来源线索；该材料可核对时间与制度称谓，但有关执行范围、群体经验、军数或后果，仍须与地方文书、物质遗存及相关政权记录互证。

\paragraph{1256年}\eventanchor{YUAN-1256-EVENT-166}{YUAN-1256-EVENT-166}　蒙哥朝加强汉地户籍、赋税和征发整顿。账本列出《元史》〈本纪〉及相关志；《元朝秘史》；《元典章》或《通制条格》相关条作为来源线索；该材料可核对时间与制度称谓，但有关执行范围、群体经验、军数或后果，仍须与地方文书、物质遗存及相关政权记录互证。

\paragraph{1256年}\eventanchor{YUAN-1256-REGNAL-YEAR}{YUAN-1256-REGNAL-YEAR}　蒙古以蒙哥在位第6年作为诏令、赋役与外交文书的纪年框架。账本列出《元史》〈本纪〉及相关志；《元朝秘史》；《元典章》或《通制条格》相关条作为来源线索；该材料可核对时间与制度称谓，但有关执行范围、群体经验、军数或后果，仍须与地方文书、物质遗存及相关政权记录互证。

\paragraph{1257年}\eventanchor{YUAN-1257-REGNAL-YEAR}{YUAN-1257-REGNAL-YEAR}　蒙古以蒙哥在位第7年作为诏令、赋役与外交文书的纪年框架。账本列出《元史》〈本纪〉及相关志；《元朝秘史》；《元典章》或《通制条格》相关条作为来源线索；该材料可核对时间与制度称谓，但有关执行范围、群体经验、军数或后果，仍须与地方文书、物质遗存及相关政权记录互证。

\paragraph{1258年}\eventanchor{YUAN-1258-REGNAL-YEAR}{YUAN-1258-REGNAL-YEAR}　蒙古以蒙哥在位第8年作为诏令、赋役与外交文书的纪年框架。账本列出《元史》〈本纪〉及相关志；《元朝秘史》；《元典章》或《通制条格》相关条作为来源线索；该材料可核对时间与制度称谓，但有关执行范围、群体经验、军数或后果，仍须与地方文书、物质遗存及相关政权记录互证。

\paragraph{1259年}\eventanchor{YUAN-1259-EVENT-167}{YUAN-1259-EVENT-167}　蒙哥在攻宋期间去世。账本列出《元史》〈本纪〉及相关志；《元朝秘史》；《元典章》或《通制条格》相关条作为来源线索；该材料可核对时间与制度称谓，但有关执行范围、群体经验、军数或后果，仍须与地方文书、物质遗存及相关政权记录互证。

\paragraph{1259年}\eventanchor{YUAN-1259-REGNAL-YEAR}{YUAN-1259-REGNAL-YEAR}　蒙古以蒙哥在位第9年作为诏令、赋役与外交文书的纪年框架。账本列出《元史》〈本纪〉及相关志；《元朝秘史》；《元典章》或《通制条格》相关条作为来源线索；该材料可核对时间与制度称谓，但有关执行范围、群体经验、军数或后果，仍须与地方文书、物质遗存及相关政权记录互证。

\paragraph{1260年}\eventanchor{YUAN-1260-EVENT-168}{YUAN-1260-EVENT-168}　忽必烈即位，与阿里不哥展开争位冲突。账本列出《元史》〈本纪〉及相关志；《元朝秘史》；《元典章》或《通制条格》相关条作为来源线索；该材料可核对时间与制度称谓，但有关执行范围、群体经验、军数或后果，仍须与地方文书、物质遗存及相关政权记录互证。

\paragraph{1260年}\eventanchor{YUAN-1260-REGNAL-YEAR}{YUAN-1260-REGNAL-YEAR}　蒙古／元以忽必烈在位第1年作为诏令、赋役与外交文书的纪年框架。账本列出《元史》〈本纪〉及相关志；《元朝秘史》；《元典章》或《通制条格》相关条作为来源线索；该材料可核对时间与制度称谓，但有关执行范围、群体经验、军数或后果，仍须与地方文书、物质遗存及相关政权记录互证。

\paragraph{1261年}\eventanchor{YUAN-1261-EVENT-169}{YUAN-1261-EVENT-169}　李璮在山东反叛后被镇压。账本列出《元史》〈本纪〉及相关志；《元朝秘史》；《元典章》或《通制条格》相关条作为来源线索；该材料可核对时间与制度称谓，但有关执行范围、群体经验、军数或后果，仍须与地方文书、物质遗存及相关政权记录互证。

\paragraph{1261年}\eventanchor{YUAN-1261-REGNAL-YEAR}{YUAN-1261-REGNAL-YEAR}　蒙古／元以忽必烈在位第2年作为诏令、赋役与外交文书的纪年框架。账本列出《元史》〈本纪〉及相关志；《元朝秘史》；《元典章》或《通制条格》相关条作为来源线索；该材料可核对时间与制度称谓，但有关执行范围、群体经验、军数或后果，仍须与地方文书、物质遗存及相关政权记录互证。

\paragraph{1262年}\eventanchor{YUAN-1262-REGNAL-YEAR}{YUAN-1262-REGNAL-YEAR}　蒙古／元以忽必烈在位第3年作为诏令、赋役与外交文书的纪年框架。账本列出《元史》〈本纪〉及相关志；《元朝秘史》；《元典章》或《通制条格》相关条作为来源线索；该材料可核对时间与制度称谓，但有关执行范围、群体经验、军数或后果，仍须与地方文书、物质遗存及相关政权记录互证。

\paragraph{1263年}\eventanchor{YUAN-1263-REGNAL-YEAR}{YUAN-1263-REGNAL-YEAR}　蒙古／元以忽必烈在位第4年作为诏令、赋役与外交文书的纪年框架。账本列出《元史》〈本纪〉及相关志；《元朝秘史》；《元典章》或《通制条格》相关条作为来源线索；该材料可核对时间与制度称谓，但有关执行范围、群体经验、军数或后果，仍须与地方文书、物质遗存及相关政权记录互证。

\paragraph{1264年}\eventanchor{YUAN-1264-EVENT-170}{YUAN-1264-EVENT-170}　阿里不哥失败，忽必烈结束主要继承战争。账本列出《元史》〈本纪〉及相关志；《元朝秘史》；《元典章》或《通制条格》相关条作为来源线索；该材料可核对时间与制度称谓，但有关执行范围、群体经验、军数或后果，仍须与地方文书、物质遗存及相关政权记录互证。

\paragraph{1264年}\eventanchor{YUAN-1264-REGNAL-YEAR}{YUAN-1264-REGNAL-YEAR}　蒙古／元以忽必烈在位第5年作为诏令、赋役与外交文书的纪年框架。账本列出《元史》〈本纪〉及相关志；《元朝秘史》；《元典章》或《通制条格》相关条作为来源线索；该材料可核对时间与制度称谓，但有关执行范围、群体经验、军数或后果，仍须与地方文书、物质遗存及相关政权记录互证。

\paragraph{1265年}\eventanchor{YUAN-1265-REGNAL-YEAR}{YUAN-1265-REGNAL-YEAR}　蒙古／元以忽必烈在位第6年作为诏令、赋役与外交文书的纪年框架。账本列出《元史》〈本纪〉及相关志；《元朝秘史》；《元典章》或《通制条格》相关条作为来源线索；该材料可核对时间与制度称谓，但有关执行范围、群体经验、军数或后果，仍须与地方文书、物质遗存及相关政权记录互证。

\paragraph{1266年}\eventanchor{YUAN-1266-REGNAL-YEAR}{YUAN-1266-REGNAL-YEAR}　蒙古／元以忽必烈在位第7年作为诏令、赋役与外交文书的纪年框架。账本列出《元史》〈本纪〉及相关志；《元朝秘史》；《元典章》或《通制条格》相关条作为来源线索；该材料可核对时间与制度称谓，但有关执行范围、群体经验、军数或后果，仍须与地方文书、物质遗存及相关政权记录互证。

\paragraph{1267年}\eventanchor{YUAN-1267-EVENT-171}{YUAN-1267-EVENT-171}　忽必烈营建大都。账本列出《元史》〈本纪〉及相关志；《元朝秘史》；《元典章》或《通制条格》相关条作为来源线索；该材料可核对时间与制度称谓，但有关执行范围、群体经验、军数或后果，仍须与地方文书、物质遗存及相关政权记录互证。

\paragraph{1267年}\eventanchor{YUAN-1267-REGNAL-YEAR}{YUAN-1267-REGNAL-YEAR}　蒙古／元以忽必烈在位第8年作为诏令、赋役与外交文书的纪年框架。账本列出《元史》〈本纪〉及相关志；《元朝秘史》；《元典章》或《通制条格》相关条作为来源线索；该材料可核对时间与制度称谓，但有关执行范围、群体经验、军数或后果，仍须与地方文书、物质遗存及相关政权记录互证。

\paragraph{1268年}\eventanchor{YUAN-1268-EVENT-172}{YUAN-1268-EVENT-172}　蒙古军围攻襄阳、樊城。账本列出《元史》〈本纪〉及相关志；《元朝秘史》；《元典章》或《通制条格》相关条作为来源线索；该材料可核对时间与制度称谓，但有关执行范围、群体经验、军数或后果，仍须与地方文书、物质遗存及相关政权记录互证。

\paragraph{1268年}\eventanchor{YUAN-1268-REGNAL-YEAR}{YUAN-1268-REGNAL-YEAR}　蒙古／元以忽必烈在位第9年作为诏令、赋役与外交文书的纪年框架。账本列出《元史》〈本纪〉及相关志；《元朝秘史》；《元典章》或《通制条格》相关条作为来源线索；该材料可核对时间与制度称谓，但有关执行范围、群体经验、军数或后果，仍须与地方文书、物质遗存及相关政权记录互证。

